\scanpage{323}
$(\unsure{1})\Longleftrightarrow(\unsure{3})$ déjà vu.

$(2)\Longleftrightarrow(\unsure{3})$ La donnée d’un hom de foncteurs cov
\[
F_A\overset{u}{\longrightarrow}\deleted{\unsure{\mathrm{End}}}\ G=F_A\circ\Gamma^p
\]
équivaut à la donnée d’un élt
\[
u_p\in\unsure{G_A|G}(A)=\operatorname{Hom}(A,\Gamma^p(A))
\]
Il faut exprimer les cond. de (2).
\inserted{Les \unsure{égalités} $f^{(0)}=1$, $f^{(1)}=0$ expriment $\unsure{u_0}=1$, $\unsure{u_1}=\mathrm{id}$.}
La relation $(f+g)^{(p)}=\sum f^{(i)}g^{(j)}$ : il suffit de voir le cas universel
\[
M=A\times A,\quad f=\mathrm{inj}_1=\alpha,\quad g=\mathrm{inj}_2=\beta,
\]
\[
\unsure{(f+g)^\beta|(f+g)}=\Delta:A\longrightarrow A\times A=M
\]
\[
(f+g)^{\unsure{n|p}}=\Gamma^p(\Delta)\circ u_p\quad\deleted{\unsure{\Gamma^p(A)}}\longrightarrow\Gamma^p(M)
\]
\[
\diagram{\begin{tikzcd}[column sep=large,row sep=large]
A\arrow[r,"u_p"]\arrow[d,"\Delta_A^{\unsure{(p)}}"']&\Gamma^p(A)\arrow[r,"\Gamma^p(\Delta_A)"]\arrow[dr,"\Delta_{\Gamma^*A}"]&\Gamma^p(A\unsure{\oplus|\times}A)\arrow[d,"\simeq"]\\
\underbrace{A\times-\times A}_{p}\arrow[rr]&{}&\displaystyle\sum_{i+j=p}\Gamma^{\deleted{\unsure{p}}i}A\otimes\Gamma^jA
\end{tikzcd}}
\]
\margin{\unsure{car} $A\longrightarrow\Gamma^i(A)\otimes\Gamma^j(A)$}
\figure{\begin{tikzpicture}\draw (0,.8)..controls(-.3,.6)and(-.2,.3)..(0,0);\draw(2,.8)..controls(2.3,.4)and(2.3,0)..(2,-.2);\end{tikzpicture}}
On trouve $\Delta(u_p)=\displaystyle\sum_{i+j=p}\unsure{u^i|u_i}\otimes u_j$
\margin{324}
